\chapter{Hamilton--Jacobi theory and transport equations}\label{ch:hj}
\providecommand{\osc}[1]{\mathcal{I}[#1]}

Elliptic operators have no singular solutions: if $P$ is elliptic and
$Pu\in\CI$ then $u\in\CI$ (\cref{ch:manifolds}).  For non-elliptic operators,
such as the wave operator, singular solutions do exist, and the simplest of them
are conormal distributions.  In this chapter we find out when, and how, one can
solve $Pu\in\CI$ with $u\in\Ic{m}(X,Y;\Omegah)$ conormal to a hypersurface $Y$.
The principal symbol of $Pu$ vanishes exactly when $Y$ is \emph{characteristic};
characteristic hypersurfaces are found by solving the eikonal equation
$p(x,dh)=0$, which we do by the Hamilton--Jacobi method, i.e.\ by symplectic
geometry on $T^*X$.  The next term in the symbol of $Pu$ is then given by a
first-order differential operator along the Hamilton vector field of $p$ on
$N^*Y$, and requiring it to vanish leads to the \emph{transport equations}.
Solving these iteratively produces solutions of $Pu\in\CI$ with prescribed
restriction to a transversal hypersurface (``travelling waves''); this is used in
\cref{ch:wave} to construct a parametrix for the wave equation.

We use the conormal distributions of \cref{ch:conormal} (whose notation and
normalization of orders we follow) and the pseudodifferential operators on
manifolds of \cref{ch:manifolds}.  Throughout, $X$ is a manifold of dimension
$n$, $\pi:T^*X\to X$ is the projection, and operators act on half-densities.

\section{Pseudodifferential operators on conormal distributions}\label{hj:sec:psido}

\begin{theorem}\label{hj:thm:psido-conormal}
  Let $Y\subset X$ be a closed embedded submanifold and
  $A\in\Psido{m'}(X;\Omegah)$ properly supported.  Then
  \begin{equation*}
    A:\Ic{m}(X,Y;\Omegah)\longrightarrow\Ic{m+m'}(X,Y;\Omegah),\qquad
    \sigma_{m+m'}(Au)=\sigma_{m'}(A)|_{N^*Y}\,\sigma_m(u),
  \end{equation*}
  and $A$ preserves the classical subspaces if it has a polyhomogeneous symbol.
  The product on the right makes sense since $\sigma_{m'}(A)\in
  S^{m'-[1]}(T^*X)$ restricts to a class in $S^{m'-[1]}(N^*Y)$.
\end{theorem}
\begin{proof}
  By a partition of unity we may assume $u$ has compact support in an adapted
  chart.  Since $A$ is properly supported and its kernel is smooth off the
  diagonal, $\phi Au$ is smooth if $\phi\in\CI$ vanishes near $\supp u$, and for
  $\phi\in\CIc$ of the chart, $\phi A\chi$ (with $\chi=1$ near $\supp u$) is, in
  the chart, $\Op(b)$ plus a smoothing operator, with $b\in S^{m'}(\R^n;\R^n)$
  compactly supported in $z$ (\cref{ch:manifolds}).  The result therefore follows
  from the local computation of \cref{ch:conormal}: if $u=(\osc{a}+f)|dz|^{1/2}$
  then $\Op(b)u$ has full symbol congruent to $b(0,y,\xi,0)a(y,\xi)$ modulo
  $S^{M+m'-1}$, and $b(0,y,\xi,0)$ represents $\sigma_{m'}(A)$ at
  $(0,y;\xi,0)\in N^*Y$.  Here and below $\osc{a}(x,y)=\int e^{ix\cdot\xi}
  a(y,\xi)\,\dbar\xi$ in adapted coordinates $(x,y)\in\R^k\times\R^d$, and
  $M=m+\frac n4-\frac k2$.
\end{proof}

In particular, if $A$ is elliptic and $\sigma_m(u)$ is not in
$S^{m+\frac n4-1}$, then $Au\notin\CI$: there are no conormal solutions of
$Au\in\CI$.  Suppose instead that $P\in\Psido{m'}(X;\Omegah)$ is not elliptic,
with homogeneous principal symbol $p$, and that we look for
$u\in\Ic{m}(X,Y;\Omegah)$ which is genuinely singular, say with
$\sigma_m(u)$ elliptic (bounded below by $c\ang{\xi}^{m+\frac n4}|dy\,d\xi|^{1/2}$
for large $\xi$), and $Pu\in\CI$.  By \cref{hj:thm:psido-conormal},
$p|_{N^*Y}\sigma_m(u)=0$ modulo lower order, so $p$ must vanish on $N^*Y$.

\begin{definition}\label{hj:def:characteristic}
  Let $P\in\Psido{m'}(X;\Omegah)$ have principal symbol $p$, homogeneous of degree
  $m'$ on $T^*X\setminus0$, and characteristic variety
  $\Sigma(P)=\{p=0\}\subset T^*X\setminus0$.  A closed embedded submanifold $Y$ is
  \emph{characteristic} for $P$ if $N^*Y\setminus0\subset\Sigma(P)$, and
  characteristic on an open cone $\Gamma\subset N^*Y\setminus0$ if
  $\Gamma\subset\Sigma(P)$.
\end{definition}

\begin{corollary}\label{hj:cor:char-lower}
  If $Y$ is characteristic for $P$ on the open cone $\Gamma$, and
  $u\in\Ic{m}(X,Y;\Omegah)$ has $\WF(u)\subset\Gamma$, then
  $Pu\in\Ic{m+m'-1}(X,Y;\Omegah)$.
\end{corollary}
\begin{proof}
  By \cref{ch:conormal} (the characterization of the wavefront set of a
  conormal distribution by its full symbol), locally the full symbol $a$ of $u$
  is of order $-\infty$ outside any conic neighbourhood of the closed cone
  $\WF(u)$; choosing a cut-off $\psi$ homogeneous of degree $0$, equal to one near
  $\WF(u)$ and supported in $\Gamma$, $a=\psi a$ modulo $S^{-\infty}$, and
  $p\psi=0$ for large $\xi$.  Hence $\sigma_{m+m'}(Pu)=p|_{N^*Y}\sigma_m(u)=0$ and
  $Pu\in\Ic{m+m'-1}$ by the exact sequence of \cref{ch:conormal}.
\end{proof}

So for characteristic $Y$ we only gain one order; we return to this in
\cref{hj:sec:transport}.  First we note that characteristic hypersurfaces and
singular solutions do exist.

\begin{example}[Plane waves]\label{hj:ex:plane}
  Let $P=D_t^2-\sum_{j=1}^nD_{x_j}^2=D_t^2-\Delta$ be the wave operator on
  $\R^{n+1}=\R_t\times\R^n_x$, acting on half-densities via the trivialization
  $|dt\,dx|^{1/2}$, so $p(t,x;\tau,\xi)=\tau^2-|\xi|^2$.  A hypersurface
  $H=\{h(t,x)=0\}$ with $dh\ne0$ on $H$ has $N^*H=\{(t,x;\lambda\,dh(t,x));\
  h(t,x)=0,\ \lambda\ne0\}$, so $H$ is characteristic if and only if $h$ satisfies
  the \emph{eikonal equation}
  \begin{equation*}
    \Big(\frac{\pa h}{\pa t}\Big)^2=|d_xh|^2\quad\text{on } \{h=0\}.
  \end{equation*}
  Linear solutions are $h=at+b\cdot x$ with $a^2=|b|^2$; up to scaling these are
  the planes $t=\omega\cdot x$, $\omega\in\Sphere^{n-1}$.  For
  $v\in\Ic{m_0}(\R,\{0\})$ (a conormal distribution at a point, in one variable),
  \begin{equation*}
    u=v(t-\omega\cdot x)\,|dt\,dx|^{1/2}\in\Ic{m_0-\frac n4}(\R^{n+1},H;\Omegah),
    \qquad Pu=(1-|\omega|^2)v''(t-\omega\cdot x)|dt\,dx|^{1/2}=0.
  \end{equation*}
  (The order follows from the linear change of coordinates $s=t-\omega\cdot x$,
  $x'=x$, of determinant one, and the relation $M=m+\frac{n+1}4-\frac12$ between
  the order of $u$ and that of its symbol, which is $m_0-\frac14$.)  For instance
  $\delta(t-\omega\cdot x)$ is a travelling wave.  We shall find all characteristic
  hypersurfaces locally, by the Hamilton--Jacobi method.
\end{example}

\section{Symplectic geometry}\label{hj:sec:symplectic}

\subsection{The canonical forms}
The cotangent bundle carries a canonical one-form $\alpha$, defined by
\begin{equation*}
  \alpha_{(x,\zeta)}(v)=\zeta(\pi_*v),\qquad v\in T_{(x,\zeta)}(T^*X).
\end{equation*}
If $x_1,\dots,x_n$ are local coordinates on $X$ and $\zeta=\sum_j\xi_jdx_j$, then
$(x,\xi)$ are coordinates on $T^*X$ and $\alpha=\sum_j\xi_j\,dx_j$.  The
\emph{symplectic form} is
\begin{equation*}
  \omega=d\alpha=\sum_{j=1}^nd\xi_j\wedge dx_j .
\end{equation*}
A section $\theta:X\to T^*X$ is the same thing as a one-form on $X$, and the
defining property of $\alpha$ gives $\theta^*\alpha=\theta$, hence
$\theta^*\omega=d\theta$.

A \emph{symplectic manifold} is a manifold $S$ with a closed two-form $\omega$
which is non-degenerate at each point; then $\dim S=2n$ is even and
$\omega^{\wedge n}\neq0$.  A \emph{symplectic vector space} is a real vector space
$E$ with a non-degenerate antisymmetric bilinear form $\omega$; non-degeneracy
means that $\tilde\omega:E\to E'$, $\tilde\omega(v)(w)=\omega(v,w)$, is an
isomorphism.  Each tangent space $T_zS$ of a symplectic manifold is a symplectic
vector space.  For a subspace $W\subset E$ set
$W^\omega=\{v\in E;\ \omega(v,w)=0\ \forall\ w\in W\}$.  A subspace is
\emph{Lagrangian} if $\omega|_{W\times W}=0$ and $2\dim W=\dim E$.

\begin{lemma}\label{hj:lem:lagrangian-linear}
  If $W$ is a subspace of a symplectic vector space $E$ then $\dim W+\dim
  W^\omega=\dim E$, and $W$ is Lagrangian if and only if $W^\omega=W$.
\end{lemma}
\begin{proof}
  Let $W^0\subset E'$ be the annihilator of $W$, of dimension $\dim E-\dim W$.
  By definition $v\in W^\omega$ if and only if $\tilde\omega(v)\in W^0$, so
  $\tilde\omega$ restricts to an isomorphism $W^\omega\to W^0$ and the dimension
  formula follows.  If $W$ is Lagrangian then $W\subset W^\omega$ and the two
  have the same dimension, so they are equal.  Conversely if $W=W^\omega$ then
  $\omega$ vanishes on $W$, and $2\dim W=\dim E$ by the dimension formula.
\end{proof}

A submanifold $\Lambda$ of a symplectic manifold is \emph{Lagrangian} if
$T_\lambda\Lambda$ is a Lagrangian subspace of $T_\lambda S$ for every
$\lambda\in\Lambda$, i.e.\ $\dim\Lambda=\frac12\dim S$ and the pull-back of
$\omega$ to $\Lambda$ vanishes.

\begin{lemma}\label{hj:lem:examples}
  The fibres $T^*_xX$ and, for a submanifold $Y\subset X$, the conormal bundle
  $N^*Y$ are Lagrangian submanifolds of $T^*X$.
\end{lemma}
\begin{proof}
  On a fibre $\pi_*v=0$ for tangent vectors, so $\alpha$ pulls back to zero, hence
  so does $\omega=d\alpha$.  On $N^*Y$, if $v$ is tangent at $(y,\zeta)$ then
  $\pi_*v\in T_yY$ and $\zeta$ annihilates $T_yY$, so again $\alpha$ pulls back to
  zero.  Both have dimension $n$.
\end{proof}

\begin{lemma}\label{hj:lem:graph}
  Let $\Lambda\subset T^*X$ be a submanifold of dimension $n$ such that
  $\pi|_\Lambda$ is a local diffeomorphism near $\bar\lambda\in\Lambda$.  Then
  near $\bar\lambda$, $\Lambda=\{(x,dh(x))\}$ for some $h$ defined near
  $\pi(\bar\lambda)$, unique up to an additive constant, if and only if $\Lambda$
  is Lagrangian near $\bar\lambda$.
\end{lemma}
\begin{proof}
  Near $\bar\lambda$, $\Lambda$ is the image of a section $\theta$ of $T^*X$ over a
  neighbourhood $U$ of $\pi(\bar\lambda)$, and $\theta:U\to\Lambda$ is a
  diffeomorphism with $\theta^*\omega=d\theta$.  So $\Lambda$ is Lagrangian if and
  only if the one-form $\theta$ is closed, which on a ball is equivalent to
  $\theta=dh$.
\end{proof}

\subsection{Hamilton vector fields}
On a symplectic manifold the isomorphism $\tilde\omega$ associates to each
function $f$ its \emph{Hamilton vector field} $H_f$, defined by
\begin{equation*}
  \omega(v,H_f)=df(v)\quad\forall\ v\in T_zS.
\end{equation*}
On $T^*X$, in canonical coordinates,
\begin{equation}\label{hj:eq:Hamilton}
  H_f=\sum_{j=1}^n\Big(\frac{\pa f}{\pa\xi_j}\frac{\pa}{\pa x_j}
  -\frac{\pa f}{\pa x_j}\frac{\pa}{\pa\xi_j}\Big).
\end{equation}
By antisymmetry $H_ff=df(H_f)=\omega(H_f,H_f)=0$: $f$ is constant along the
integral curves of $H_f$.  The flow $\exp(sH_f)$ preserves $\omega$, since by
Cartan's formula $\mathcal L_{H_f}\omega=d\iota_{H_f}\omega+\iota_{H_f}d\omega
=-d\,df=0$.  If $f$ is real, $H_f$ is a real vector field; if $f$ is homogeneous of
degree $\mu$ in the fibre variables then by \eqref{hj:eq:Hamilton} $H_f$ is
homogeneous of degree $\mu-1$ (its $\pa_x$-components are homogeneous of degree
$\mu-1$ and its $\pa_\xi$-components of degree $\mu$), and its flow commutes with
the dilations of the fibres up to reparametrization.  The condition
$d_{\mathrm{fib}}f\ne0$, i.e.\ $(\pa f/\pa\xi_1,\dots,\pa f/\pa\xi_n)\ne0$,
means precisely that $\pi_*H_f\ne0$, i.e.\ $H_f$ is transversal to the fibres
of $T^*X$.

\begin{lemma}\label{hj:lem:tangent}
  If $\Lambda\subset S$ is Lagrangian and $f\in\CI(S)$ vanishes on $\Lambda$,
  then $H_f$ is tangent to $\Lambda$.
\end{lemma}
\begin{proof}
  For $v\in T_\lambda\Lambda$, $\omega(v,H_f)=df(v)=0$; so
  $H_f(\lambda)\in(T_\lambda\Lambda)^\omega=T_\lambda\Lambda$ by
  \cref{hj:lem:lagrangian-linear}.
\end{proof}

In particular, if $Y$ is characteristic for $P$ on $\Gamma\subset N^*Y$, then
$H_p$ is tangent to $\Gamma$: a characteristic conormal bundle is a union of
integral curves of $H_p$ (bicharacteristics).  We use this to construct
characteristic hypersurfaces.

\section{Hamilton--Jacobi theory}\label{hj:sec:HJ}

Let $p\in\CI(W)$ be a real-valued function on an open set $W\subset T^*X$.  We
seek solutions of the \emph{eikonal equation}
\begin{equation}\label{hj:eq:eikonal}
  p(x,dh(x))=0
\end{equation}
in a neighbourhood of a point $\bar x$, with $h$ prescribed on a hypersurface
through $\bar x$.  Note that \eqref{hj:eq:eikonal} is required on an open set,
not only on $\{h=0\}$.  By \cref{hj:lem:graph}, solving it amounts to finding a
Lagrangian submanifold $\Lambda\subset p^{-1}(0)$ which projects
diffeomorphically to a neighbourhood of $\bar x$; by \cref{hj:lem:tangent} such
a $\Lambda$ must be a union of integral curves of $H_p$.

\begin{theorem}[Hamilton--Jacobi]\label{hj:thm:HJ}
  Let $p\in\CI(W)$ be real, $\bar z\in W$ with $p(\bar z)=0$, $\bar x=\pi(\bar
  z)$, and let $H=\{t=0\}$ be a hypersurface through $\bar x$, $t\in\CI(X)$ real
  with $dt\neq0$, such that
  \begin{equation}\label{hj:eq:transversal}
    H_p(\pi^*t)(\bar z)\neq0.
  \end{equation}
  Let $f\in\CI(H)$ satisfy $df(\bar x)=\bar z|_{T_{\bar x}H}$.  Then there are a
  neighbourhood $O$ of $\bar x$ and $h\in\CI(O)$ such that
  \begin{equation*}
    p(x,dh(x))=0\ \text{ in } O,\qquad h|_{H\cap O}=f|_{H\cap O},\qquad
    dh(\bar x)=\bar z,
  \end{equation*}
  and any other such solution coincides with $h$ near $\bar x$.  Moreover the
  graph of $dh$ is the union of the integral curves of $H_p$ through the
  submanifold $\Lambda_0$ defined in \eqref{hj:eq:Lambda0} below.
\end{theorem}
\begin{proof}
  Choose coordinates $(t,x')$ near $\bar x=0$, $x'\in\R^{n-1}$, with dual
  coordinates $(\tau,\xi')$; then $H=\{t=0\}$, $\bar z=(0,0;\bar\tau,\bar\xi')$
  with $\bar\xi'=d_{x'}f(0)$, and \eqref{hj:eq:transversal} says
  $\pa_\tau p(\bar z)\neq0$.

  \emph{Initial manifold.}  By the implicit function theorem there is a unique
  smooth function $\tau(x')$, defined near $0$, with $\tau(0)=\bar\tau$ and
  $p(0,x';\tau(x'),d_{x'}f(x'))=0$.  Set
  \begin{equation}\label{hj:eq:Lambda0}
    \Lambda_0=\{(0,x';\tau(x'),d_{x'}f(x'))\},
  \end{equation}
  an $(n-1)$-dimensional submanifold of $p^{-1}(0)$.  It consists exactly of the
  points of $p^{-1}(0)$ near $\bar z$ lying over $H$ whose restriction to $TH$ is
  $df$.  The pull-back of $\omega$ to $\Lambda_0$ is $d\tau\wedge dt+\sum
  d\xi'_j\wedge dx'_j$ restricted, which is $0+d(d_{x'}f)=0$ since $dt=0$ on
  $\Lambda_0$.

  \emph{Flow-out.}  Let $\Phi_s=\exp(sH_p)$ and $\Psi(s,x')=\Phi_s(0,x';\tau(x'),
  d_{x'}f(x'))$ for $|s|<\epsilon$ and $x'$ near $0$.  At $(0,0)$, $\pa_s\Psi=
  H_p(\bar z)$, whose projection to $X$ has $\pa_t$-component $\pa_\tau
  p(\bar z)\neq0$, while $\pa_{x'_j}\Psi$ projects to $\pa_{x'_j}$.  Hence
  $\pi\circ\Psi$ has invertible differential at $(0,0)$, so, shrinking the
  domain, $\Psi$ is an embedding onto an $n$-dimensional submanifold $\Lambda$
  and $\pi|_\Lambda$ is a diffeomorphism onto a neighbourhood $O$ of $\bar x$.
  Since $H_pp=0$, $\Lambda\subset p^{-1}(0)$.

  \emph{$\Lambda$ is Lagrangian.}  At a point $z\in\Lambda_0$,
  $T_z\Lambda=T_z\Lambda_0\oplus\R H_p(z)$; $\omega$ vanishes on
  $T_z\Lambda_0\times T_z\Lambda_0$, $\omega(H_p,H_p)=0$, and for $v\in
  T_z\Lambda_0$, $\omega(v,H_p)=dp(v)=0$ because $\Lambda_0\subset p^{-1}(0)$.
  So $T_z\Lambda$ is Lagrangian.  Since $\Psi(s+s',x')=\Phi_s\Psi(s',x')$, the
  differential of $\Phi_s$ maps $T_z\Lambda$ onto $T_{\Phi_s(z)}\Lambda$, and
  $\Phi_s$ preserves $\omega$; so every tangent space of $\Lambda$ is Lagrangian.

  \emph{Conclusion.}  By \cref{hj:lem:graph}, $\Lambda$ is the graph of $dh$ for
  some $h\in\CI(O)$, and $p(x,dh)=0$.  The points of $\Lambda$ over $H$ are those
  of $\Lambda_0$ (as $\pi|_\Lambda$ is injective and $\pi(\Lambda_0)=H\cap O$), so
  $d(h|_H)=d_{x'}f$; adjusting the constant, $h|_H=f$, and $dh(\bar x)=\bar z$
  since $\bar z\in\Lambda_0$.

  \emph{Uniqueness.}  If $h'$ is another solution, its graph $\Lambda'$ is
  Lagrangian and contained in $p^{-1}(0)$, so $H_p$ is tangent to $\Lambda'$ by
  \cref{hj:lem:tangent}.  The points $(0,x';\pa_th'(0,x'),d_{x'}f(x'))$ of
  $\Lambda'$ over $H$ satisfy $p=0$ and are close to $\bar z$, so by the uniqueness
  in the implicit function theorem $\pa_th'(0,x')=\tau(x')$, i.e.\
  $\Lambda_0\subset\Lambda'$.  The integral curves of $H_p$ starting on
  $\Lambda_0$ therefore stay in $\Lambda'$, so $\Lambda\subset\Lambda'$ near
  $\bar z$.  Both are graphs over a neighbourhood of $\bar x$, hence $dh=dh'$
  there, and $h=h'$ since they agree on $H$.
\end{proof}

Note that \eqref{hj:eq:transversal} implies $d_{\mathrm{fib}}p(\bar z)\ne0$;
conversely, if $d_{\mathrm{fib}}p(\bar z)\neq0$ then $\pi_*H_p(\bar z)\neq0$ and
hypersurfaces $H$ satisfying \eqref{hj:eq:transversal} exist.

For homogeneous $p$ we obtain characteristic hypersurfaces with prescribed
intersection with $H$.  If $G\subset H$ is a hypersurface in $H$ and $\bar
x\in G$, write $N^*_{\bar x}G\subset T^*_{\bar x}H$ for its conormal space in
$H$; and for $\bar z\in T^*_{\bar x}X$ write $\bar z_H=\bar z|_{T_{\bar x}H}\in
T^*_{\bar x}H$.

\begin{corollary}\label{hj:cor:characteristic}
  Let $p\in\CI(T^*X\setminus0)$ be real and homogeneous of degree $m'$, $\bar z\in
  T^*X\setminus0$ with $p(\bar z)=0$, $H=\{t=0\}$ a hypersurface through $\bar
  x=\pi(\bar z)$ satisfying \eqref{hj:eq:transversal}, and $G\subset H$ a
  hypersurface of $H$ through $\bar x$ with $\bar z_H\in N^*_{\bar x}G\setminus0$.
  Then there is a hypersurface $\tilde G$ of $X$, defined near $\bar x$, with
  $\tilde G\cap H=G$ near $\bar x$, $\bar z\in N^*\tilde G$, and
  \begin{equation*}
    \Gamma=\{(x,\lambda\,d\tilde g(x));\ x\in\tilde G,\ \lambda>0\}\subset\Sigma(p),
  \end{equation*}
  where $\tilde g$ is a defining function of $\tilde G$ with $d\tilde g(\bar
  x)=\bar z$; that is, $\tilde G$ is characteristic on the half $\Gamma$ of
  $N^*\tilde G\setminus0$ containing $\bar z$.  Moreover $\tilde G$ is unique near
  $\bar x$.
\end{corollary}
\begin{proof}
  Write $G=\{g_0=0\}$ near $\bar x$ with $g_0\in\CI(H)$, $dg_0\ne0$; since
  $\bar z_H$ is a nonzero multiple of $dg_0(\bar x)$ we may rescale $g_0$ so that
  $dg_0(\bar x)=\bar z_H$.  \Cref{hj:thm:HJ} with $f=g_0$ gives $\tilde g$ with
  $p(x,d\tilde g)=0$, $\tilde g|_H=g_0$ and $d\tilde g(\bar x)=\bar z\ne0$.  Then
  $\tilde G=\{\tilde g=0\}$ is a hypersurface near $\bar x$, $\tilde G\cap
  H=\{g_0=0\}=G$, and $p(x,\lambda d\tilde g(x))=\lambda^{m'}p(x,d\tilde g(x))=0$
  for $\lambda>0$.

  For uniqueness, let $\tilde G'$ be another such hypersurface, with half-conormal
  bundle $\Gamma'\subset\Sigma(p)$ containing $\bar z$.  Over $G$, the points of
  $\Gamma$ and of $\Gamma'$ near $\bar z$ are covectors in $\Sigma(p)$ whose
  restrictions to $TH$ lie in $N^*G$; by \eqref{hj:eq:transversal} and the
  implicit function theorem (as for $\Lambda_0$ above, now with the extra
  dilation parameter) these sets coincide near $\bar z$; call this
  $(n-1)$-dimensional manifold $\Gamma_0$.  Since $\Gamma,\Gamma'$ are Lagrangian
  (\cref{hj:lem:examples}) and contained in $p^{-1}(0)$, $H_p$ is tangent to both,
  so both contain the flow-out of $\Gamma_0$ by $H_p$.  This flow-out is
  $n$-dimensional near $\bar z$: $H_p(\bar z)$ is not tangent to $\Gamma_0$ since
  $\pi_*H_p(\bar z)\notin T_{\bar x}H\supset\pi_*T\Gamma_0$.  So it is an open
  subset of both $\Gamma$ and $\Gamma'$, and $\tilde G=\pi(\Gamma)=\pi(\Gamma')=
  \tilde G'$ near $\bar x$.
\end{proof}

\section{The transport equation}\label{hj:sec:transport}

\subsection{The symbol of $Pu$ for characteristic $Y$}
We now compute the next term in the symbol of $Pu$ when $Y$ is characteristic.
We assume that $P\in\Psido{m'}(X;\Omegah)$ is properly supported and has a
polyhomogeneous symbol: in local coordinates, acting on coefficients with
respect to $|dz|^{1/2}$, $P=\Op(b)$ modulo smoothing operators, with
\begin{equation*}
  b\sim b_{m'}+b_{m'-1}+\cdots,\qquad b_{m'-j}\text{ homogeneous of degree }
  m'-j\text{ in }|\zeta|\ge1 .
\end{equation*}
Differential operators are the main examples.  The principal symbol is
$p=b_{m'}$, and the \emph{subprincipal symbol} is
\begin{equation}\label{hj:eq:subprincipal}
  p_s=b_{m'-1}-\frac1{2i}\sum_{j=1}^n\frac{\pa^2p}{\pa z_j\pa\zeta_j}.
\end{equation}
For operators on half-densities $p_s$ is an invariantly defined function on
$T^*X\setminus0$, homogeneous of degree $m'-1$ \cite[Theorem~18.1.33]{Hormander-III};
we shall only use its restriction to $N^*Y$, whose invariance follows from the
next theorem.  Recall that the Lie derivative acts on half-densities on a
manifold with coordinates $w$ by
$\mathcal L_V(\beta|dw|^{1/2})=(V\beta+\frac12(\operatorname{div}V)\beta)|dw|^{1/2}$,
$\operatorname{div}V=\sum\pa_{w_j}V_j$.

\begin{theorem}\label{hj:thm:transport}
  Let $P$ be as above, $Y\subset X$ a closed embedded submanifold and
  $\Gamma\subset N^*Y\setminus0$ an open cone on which $p=0$, so that $H_p$ is
  tangent to $\Gamma$.  If $u\in\Ic{m}(X,Y;\Omegah)$ and $\WF(u)\subset\Gamma$,
  then $Pu\in\Ic{m+m'-1}(X,Y;\Omegah)$ and
  \begin{equation}\label{hj:eq:transport}
    \sigma_{m+m'-1}(Pu)=\big(-i\mathcal L_{H_p}+p_s\big)\sigma_m(u).
  \end{equation}
\end{theorem}
Here $\sigma_m(u)\in S^{m+\frac n4}(N^*Y;\Omegah)$ is, modulo $S^{-\infty}$,
supported in a closed cone contained in $\Gamma$, and $\mathcal L_{H_p}$ (with
$H_p$ restricted to $\Gamma$) and multiplication by $p_s$ map
$S^r(\Gamma;\Omegah)$ to $S^{r+m'-1}(\Gamma;\Omegah)$, since $H_p$ and $p_s$
are homogeneous of degree $m'-1$; so the right side of \eqref{hj:eq:transport}
is well defined in $S^{m+\frac n4+m'-1-[1]}(N^*Y;\Omegah)$.
\begin{proof}
  That $Pu\in\Ic{m+m'-1}$ is \cref{hj:cor:char-lower}.  The formula is local; in
  adapted coordinates $(x,y)$ with $u=(\osc{a}+f)|dz|^{1/2}$, the full symbol
  $a'$ of $Pu$ is given by the expansion of \cref{ch:conormal},
  \begin{equation*}
    a'\sim\sum_{\alpha,\gamma}\frac1{\alpha!\gamma!}(-D_\xi)^\alpha\pa_x^\alpha
    \Big(\pa_\eta^\gamma b(x,y,\xi,0)D_y^\gamma a(y,\xi)\Big)\Big|_{x=0}.
  \end{equation*}
  Each $\alpha$ or $\gamma$ lowers the order by one, so modulo $S^{M+m'-2}$
  only the terms with $|\alpha|+|\gamma|\le1$ and the two leading terms of $b$
  contribute.  Since $p(0,y,\xi,0)=0$ where $a$ is not of order $-\infty$,
  \begin{align*}
    a'&\equiv b_{m'-1}a+\sum_l\pa_{\eta_l}p\,D_{y_l}a
       -\sum_jD_{\xi_j}\big(\pa_{x_j}p\,a\big)\\
      &=-i\Big(\sum_l\pa_{\eta_l}p\,\pa_{y_l}a-\sum_j\pa_{x_j}p\,\pa_{\xi_j}a\Big)
       +\Big(b_{m'-1}+i\sum_j\pa_{x_j}\pa_{\xi_j}p\Big)a,
  \end{align*}
  all functions of $(y,\xi)$ evaluated at $x=0$, $\eta=0$.  On $N^*Y=\{x=0,\
  \eta=0\}$ we have $\pa_\xi p=0$ and $\pa_yp=0$ (as $p$ vanishes there), so by
  \eqref{hj:eq:Hamilton}
  \begin{equation*}
    H_p|_{N^*Y}=\sum_l\pa_{\eta_l}p\,\pa_{y_l}-\sum_j\pa_{x_j}p\,\pa_{\xi_j},
  \end{equation*}
  which is the first-order operator above.  With respect to the half-density
  $|dy\,d\xi|^{1/2}$ its divergence is $\sum_l\pa_{y_l}\pa_{\eta_l}p-\sum_j
  \pa_{\xi_j}\pa_{x_j}p$, so
  \begin{equation*}
    -iH_pa=-i\mathcal L_{H_p}a+\frac i2\Big(\sum_l\pa_{y_l}\pa_{\eta_l}p-
    \sum_j\pa_{x_j}\pa_{\xi_j}p\Big)a,
  \end{equation*}
  where $\mathcal L_{H_p}a$ denotes the coefficient of
  $\mathcal L_{H_p}(a|dy\,d\xi|^{1/2})$.  Substituting,
  \begin{equation*}
    a'\equiv-i\mathcal L_{H_p}a+\Big(b_{m'-1}+\frac i2\sum_{j=1}^n
    \frac{\pa^2p}{\pa z_j\pa\zeta_j}\Big)a=-i\mathcal L_{H_p}a+p_sa
  \end{equation*}
  modulo $S^{M+m'-2}$, since $\frac i2=-\frac1{2i}$.  This is
  \eqref{hj:eq:transport}.
\end{proof}

Thus, to obtain $Pu\in\Ic{m+m'-2}$ we must solve the \emph{transport equation}
$(-i\mathcal L_{H_p}+p_s)\sigma_m(u)=0$, a first-order ordinary differential
equation along the integral curves of $H_p$ in $\Gamma$; the lower order terms
are then removed by solving inhomogeneous transport equations.  To set up the
initial value problem we need restriction to a transversal hypersurface.

\subsection{Restriction to a transversal hypersurface}
Let $t\in\CI(X)$ be real with $dt\neq0$ on $H=\{t=0\}$.  Restriction of
half-densities to $H$ is defined using $|dt|^{1/2}$: in coordinates $(t,w)$,
\begin{equation}\label{hj:eq:restrict-hd}
  \big(\phi(t,w)|dt\,dw|^{1/2}\big)\big|_H=\phi(0,w)|dw|^{1/2};
\end{equation}
this is independent of the coordinates $w$ (\cref{hj:ex:restriction-hd}).  Let
$\tilde G\subset X$ be a hypersurface transversal to $H$, i.e.\ $dt|_{T\tilde
G}\ne0$ at each point of $G=\tilde G\cap H$; then $G$ is a hypersurface in $H$
(\cref{hj:ex:transversal}).  Near any point of $G$ there are coordinates
$(x_1,t,y')$, $y'\in\R^{n-2}$, with $\tilde G=\{x_1=0\}$ and $H=\{t=0\}$.  The
restriction map $j:N^*\tilde G|_G\to N^*G$, $\zeta\mapsto\zeta|_{TH}$, is an
isomorphism, in coordinates $(t=0,y';\xi_1)\mapsto(y';\xi_1)$.  The set
$N^*\tilde G|_G=N^*\tilde G\cap\pi^{-1}(H)$ is the hypersurface
$\{\pi^*t=0\}$ of $N^*\tilde G$, so half-densities on $N^*\tilde G$ restrict to
it as in \eqref{hj:eq:restrict-hd}, using $|d\pi^*t|^{1/2}$.

\begin{proposition}\label{hj:prop:restriction}
  With $H$, $\tilde G$, $G$ as above, restriction to $H$ defines a map
  \begin{equation*}
    \Ic{m}(X,\tilde G;\Omegah)\ni u\longmapsto u|_H\in\Ic{m+\frac14}(H,G;\Omegah),
    \qquad \sigma_{m+\frac14}(u|_H)=(j^{-1})^*\big(\sigma_m(u)|_{\{\pi^*t=0\}}\big),
  \end{equation*}
  mapping $\Ic{m-l}$ to $\Ic{m-l+\frac14}$, and $\WF(u|_H)\subset
  j(\WF(u)\cap\pi^{-1}(H))$.
\end{proposition}
\begin{proof}
  Since $\WF(u)\subset N^*\tilde G$ and $N^*\tilde G\cap N^*H=0$ by
  transversality, the restriction (pull-back by the inclusion) of $u$ to $H$ is
  well defined \cite[Theorem~8.2.4]{Hormander-I}; for half-densities we combine it
  with \eqref{hj:eq:restrict-hd}.  In coordinates as above,
  $u=(\osc{a}+f)|dx_1\,dt\,dy'|^{1/2}$ with $a=a(t,y',\xi_1)\in S^M$,
  $M=m+\frac n4-\frac12$, and $\osc{a}$ is a smooth function of $t$ with values in
  distributions in $(x_1,y')$; so
  \begin{equation*}
    u|_H=(\osc{a(0,\cdot)}+f(\cdot,0,\cdot))|dx_1\,dy'|^{1/2}.
  \end{equation*}
  Now $H$ has dimension $n-1$ and $G$ codimension one in it, so this is an element
  of $\Ic{m'}(H,G;\Omegah)$ with $M=m'+\frac{n-1}4-\frac12$, i.e.\
  $m'=m+\frac14$.  Its symbol is $a(0,y',\xi_1)|dy'\,d\xi_1|^{1/2}$, which is the
  stated formula since $\sigma_m(u)=a|dt\,dy'\,d\xi_1|^{1/2}$.  The wavefront set
  statement follows from the characterization of $\WF$ by the full symbol
  (\cref{ch:conormal}).
\end{proof}

\subsection{Travelling waves}
We now construct conormal solutions of $Pu\in\CI$ with prescribed restriction
to a transversal hypersurface.  We first solve the transport equations.

\begin{lemma}\label{hj:lem:transport-ode}
  Under the hypotheses of \cref{hj:cor:characteristic}, use coordinates
  $(x_1,t,y')$ near $\bar x$ with $\tilde G=\{x_1=0\}$, $H=\{t=0\}$, so that
  $(t,y',\xi_1)$ are coordinates on $N^*\tilde G$ and $\Gamma=\{\xi_1>0\}$.  Let
  $q$ be a smooth function on $\Gamma$ homogeneous of degree $m'-1$.  Then there
  are $\epsilon>0$ and a neighbourhood $U$ of $\bar y'=0$ such that for every
  $c\in S^{r+m'-1}$ and $\beta_0\in S^r$ (classes of half-densities on
  $\Gamma$, respectively on $\Gamma\cap\{t=0\}$, of the stated orders, supported
  in $\xi_1\ge1$) the problem
  \begin{equation*}
    (-i\mathcal L_{H_p}+q)\beta=c\ \text{ in }\{|t|<\epsilon,\ y'\in U\},\qquad
    \beta|_{t=0}=\beta_0
  \end{equation*}
  has a unique solution $\beta\in S^r$ there (restriction to $t=0$ taken as in
  \eqref{hj:eq:restrict-hd}).  If $\beta_0$ and $c$ are of order $-\infty$ outside
  closed cones $K_0$, $K$, then $\beta$ is of order $-\infty$ outside the union of
  $K$ and the $H_p$-flow-out of $K_0$.
\end{lemma}
\begin{proof}
  Since $H_p$ is tangent to $\Gamma$ and homogeneous of degree $m'-1$, in the
  coordinates $(t,y',\xi_1)$
  \begin{equation*}
    H_p=\xi_1^{m'-1}\Big(e_0(t,y')\pa_t+\sum_le_l(t,y')\pa_{y'_l}
    +e_{n}(t,y')\xi_1\pa_{\xi_1}\Big),
  \end{equation*}
  and $e_0(0,0)\neq0$ by \eqref{hj:eq:transversal}.  Writing
  $\beta=\tilde\beta|dt\,dy'\,d\xi_1|^{1/2}$ and dividing by
  $-i\xi_1^{m'-1}e_0$, the equation becomes
  \begin{equation*}
    \Big(\pa_t+\sum_l\frac{e_l}{e_0}\pa_{y'_l}+\frac{e_n}{e_0}\xi_1\pa_{\xi_1}\Big)\tilde\beta
    +\kappa(t,y')\tilde\beta=\tilde c,
  \end{equation*}
  with $\kappa$ smooth (the divergence of $H_p$ and $q$ both being $\xi_1^{m'-1}$
  times functions of $(t,y')$) and $\tilde c=i\xi_1^{1-m'}e_0^{-1}c'$, $c'$ the
  coefficient of $c$, of the same order as $\tilde\beta$.  The vector field in
  parentheses has integral curves $s\mapsto(s,Y(s,y'_0),\xi_1E(s,y'_0))$ with
  $E>0$ smooth, the $\xi_1$-component being linear in $\xi_1$ because its
  coefficient is.  Hence the solution is given explicitly by
  \begin{equation*}
    \tilde\beta(s,Y(s,y'_0),\xi_1E(s,y'_0))=e^{-k(s,y'_0)}\Big(\tilde\beta_0(y'_0,\xi_1)
    +\int_0^se^{k(\sigma,y'_0)}\tilde c(\sigma,Y(\sigma,y'_0),\xi_1E(\sigma,y'_0))
    \,d\sigma\Big),
  \end{equation*}
  $k(s,y'_0)=\int_0^s\kappa(\sigma,Y(\sigma,y'_0))\,d\sigma$, for $|s|<\epsilon$
  and $y'_0$ near $0$.  The change of variables $(s,y'_0,\xi_1)\mapsto
  (s,Y,\xi_1E)$ is a diffeomorphism near $\{s=0,\ y'_0=0\}$ which is linear and
  positive in the fibre variable, so it preserves symbol classes; hence $\beta$ has
  order $r$, and the support statement is clear from the formula.
\end{proof}

\begin{theorem}[Travelling waves]\label{hj:thm:travelling}
  Let $P\in\Psido{m'}(X;\Omegah)$ be properly supported with polyhomogeneous
  symbol and real principal symbol $p$.  Let $\bar z\in T^*X\setminus0$ with
  $p(\bar z)=0$, $\bar x=\pi(\bar z)$, let $H=\{t=0\}$ be a hypersurface through
  $\bar x$ with $H_p(\pi^*t)(\bar z)\neq0$, let $G\subset H$ be a hypersurface of
  $H$ with $\bar z_H\in N^*_{\bar x}G\setminus0$, and let $\tilde G$ be the
  characteristic hypersurface of \cref{hj:cor:characteristic}.  Then there are a
  neighbourhood $O$ of $\bar x$ and an open conic neighbourhood $\gamma$ of
  $\bar z_H$ in $T^*H\setminus0$ such that for every
  $u_0\in\Ic{m+\frac14}(H,G;\Omegah)$ with $\WF(u_0)\subset\gamma$ there exists
  $u\in\Ic{m}(X,\tilde G;\Omegah)$ with
  \begin{equation*}
    u|_{H\cap O}=u_0|_{H\cap O},\qquad Pu\in\CI(O;\Omegah),
  \end{equation*}
  and $\WF(u)\cap T^*O\subset\Gamma$.  If $P$ and $u_0$ are classical, $u$ can be
  taken classical.
\end{theorem}
\begin{proof}
  Use the coordinates of \cref{hj:lem:transport-ode}, and choose $\epsilon$, $U$
  as there and $\gamma$ so small that, in these coordinates, $j^{-1}(\gamma)$ lies
  over $\{y'\in U\}$ and within a cone whose $H_p$-flow-out for $|t|<\epsilon$
  stays in $\Gamma$.  Let $O=\{|x_1|<\epsilon',|t|<\epsilon,\ y'\in U'\}$ with
  $U'\Subset U$, and fix $\phi\in\CIc(\{|t|<\epsilon,y'\in U\})$ equal to one
  near $\bar O$.  We construct $u_j\in\Ic{m-j}(X,\tilde G;\Omegah)$, with
  wavefront sets in $\Gamma$ and supports in $\supp\phi$, such that for
  $v_N=u_0'+\dots+u_N$ (we write $u_0'$ for the first term to avoid confusion with
  the data $u_0$)
  \begin{equation}\label{hj:eq:travelling-induction}
    Pv_N\in\Ic{m+m'-N-2}\ \text{ in }O,\qquad u_0-v_N|_H\in
    \Ic{m+\frac14-N-1}(H,G;\Omegah)\ \text{ in } H\cap O.
  \end{equation}
  For $N=-1$ (with $v_{-1}=0$) the second condition holds and the first is
  vacuous with the convention $Pv_{-1}=0$.  Given $v_{N-1}$, let
  $g=-Pv_{N-1}\in\Ic{m+m'-N-1}$ in $O$ and $e=u_0-v_{N-1}|_H\in
  \Ic{m+\frac14-N}(H,G;\Omegah)$, both with wavefront sets in the appropriate
  cones.  By \cref{hj:lem:transport-ode} (with $q=p_s|_\Gamma$) solve
  \begin{equation*}
    (-i\mathcal L_{H_p}+p_s)\beta_N=\sigma_{m+m'-N-1}(g),\qquad
    \beta_N|_{t=0}=(j)^*\sigma_{m+\frac14-N}(e),
  \end{equation*}
  with $\beta_N\in S^{m+\frac n4-N}$, of order $-\infty$ outside the flow-out
  cone, and let $u_N\in\Ic{m-N}$ be $\phi$ times a distribution with principal
  symbol $\beta_N$ (\cref{ch:conormal}).  By \cref{hj:thm:transport},
  $\sigma_{m+m'-N-1}(Pu_N)=\sigma_{m+m'-N-1}(g)$ in $O$ (where $\phi=1$ and,
  $P$ being pseudolocal, $P(\phi w)-\phi Pw$ is smooth), so $Pv_N=-g+Pu_N\in
  \Ic{m+m'-N-2}$ in $O$; and by \cref{hj:prop:restriction},
  $\sigma_{m+\frac14-N}(u_N|_H)=\sigma_{m+\frac14-N}(e)$, so $e-u_N|_H\in
  \Ic{m+\frac14-N-1}$.  This proves \eqref{hj:eq:travelling-induction}.

  Let $\tilde u\sim\sum_Nu_N$ be an asymptotic sum in $\Ic{m}(X,\tilde G;\Omegah)$,
  with support in $\supp\phi$.  Then $P\tilde u\in\bigcap_N\Ic{m+m'-N}=\CI$ in
  $O$, and $u_0-\tilde u|_H\in\bigcap_N\Ic{m+\frac14-N}(H,G)=\CI$ in $H\cap O$,
  by \cref{hj:prop:restriction}.  Finally choose $w\in\CI(X;\Omegah)$ with
  $w|_H=u_0-\tilde u|_H$ near $H\cap\bar O$ (in coordinates, $w=\chi(t)\times$ the
  smooth error, extended constantly in $t$), and set $u=\tilde u+w$.  Then
  $u|_{H\cap O}=u_0$, $Pu=P\tilde u+Pw\in\CI(O)$, and $\WF(u)\cap T^*O\subset\Gamma$.
  In the classical case all symbols above are classical.
\end{proof}

The transport equation \eqref{hj:eq:transport} for the leading symbol, and its
inhomogeneous versions for the lower order terms, are thus solved by integration
along the bicharacteristics (Duhamel's principle).  In \cref{ch:wave} this
construction, applied to the wave operator with $H=\{t=0\}$, gives a
parametrix for the Cauchy problem.

\section{Exercises}\label{hj:sec:exercises}

\begin{exercise}\label{hj:ex:transversal}
  Let $\tilde G\subset X$ be a hypersurface and $t\in\CI(X)$ real with
  $dt|_{T_q\tilde G}\neq0$ for all $q\in G=\tilde G\cap\{t=0\}$.  Show that
  $H=\{t=0\}$ is a hypersurface near $G$ and that $G=H\cap\tilde G$ is a
  submanifold of codimension one in both $H$ and $\tilde G$.
\end{exercise}

\begin{exercise}\label{hj:ex:restriction-hd}
  Assuming $dt\ne0$ on $H=\{t=0\}$, show that $dt$ gives an isomorphism of line
  bundles $\Omegah(H)\cong\Omegah_H(X)\otimes|N^*H|^{-1/2}$, where
  $\Omegah_H(X)$ is the restriction to $H$ of $\Omegah(X)$ and $|N^*H|^{-1/2}$ is
  trivialized by $|dt|^{-1/2}$, and hence that \eqref{hj:eq:restrict-hd} defines a
  coordinate-independent restriction map $\CI(X;\Omegah)\to\CI(H;\Omegah)$.
\end{exercise}

\begin{exercise}
  Prove \cref{hj:prop:restriction} directly from the coordinate formula: show that
  the restriction defined in adapted coordinates is independent of the choice of
  coordinates $(x_1,t,y')$ with $t$ fixed.
\end{exercise}

\begin{exercise}
  For the wave operator $P=D_t^2-\Delta$ on $\R^{n+1}$, $H=\{t=0\}$ and
  $f(x)=\omega\cdot x$, $|\omega|=1$, find the two solutions $h_\pm$ of the
  eikonal equation $p(x,dh)=0$ with $h|_H=f$ given by \cref{hj:thm:HJ} (one for
  each choice of $\bar z$ over $0$).  Do the same with $f(x)=|x|$ near a point
  $x\ne0$, and describe the characteristic hypersurfaces through
  $\{t=0,\ |x|=1\}$.  Show that the light cone $\{t^2=|x|^2\}$ is characteristic
  away from its vertex.
\end{exercise}

\begin{exercise}
  Show that the subprincipal symbol of $D_t^2-\Delta$ on $\R^{n+1}$, acting on
  half-densities via $|dt\,dx|^{1/2}$, vanishes, and solve the transport equation
  for the plane waves of \cref{hj:ex:plane}.  For a Riemannian metric $g$ on $X$
  compute the subprincipal symbol of $D_t^2-\Delta_g$ acting on half-densities by
  $u\mapsto(D_t^2-\Delta_g)(u\,|dg|^{-1/2})\,|dg|^{1/2}$, $|dg|$ the Riemannian
  density.
\end{exercise}

\section*{Notes}
This chapter is a rewriting of the second half of the conormal/wave kernel
chapter of the 1998 notes.  The section on Riemann metrics and quantization
there was only a heading and has been omitted.  The Hamilton--Jacobi theorem is
proved here in full (the original only sketched the strategy), the transport
equation is derived for submanifolds of any codimension with the zeroth order
term identified as the subprincipal symbol, the restriction theorem (an exercise
in the original) is proved, and the construction of travelling waves is made
precise.  For the general theory see \cite[Sections~18.1, 25.2]{Hormander-III}
and \cite{Hormander-FIO}, where conormal distributions are replaced by
Lagrangian distributions and the analogue of \eqref{hj:eq:transport} holds on
any Lagrangian submanifold contained in $\Sigma(P)$.
